% Unidad R570: introducción, líneas francesas 122--132, p. 9.
% Traducción: OpenAI Codex, GPT-6.1 Sol, esfuerzo Ultra.
Para terminar, creemos útil advertir a los lectores que, al igual que los propios autores, seguramente encontrarán alguna dificultad antes de acostumbrarse al lenguaje de los esquemas y de convencerse de que las construcciones habituales sugeridas por la intuición geométrica pueden trasladarse a este lenguaje, esencialmente de una única manera razonable.
Como en muchas partes de la matemática moderna, la intuición inicial parece alejarse cada vez más del lenguaje adecuado para expresarla con toda la precisión y generalidad deseadas.
En este caso, la dificultad psicológica se debe a la necesidad de trasladar a los objetos de una categoría ya bastante distinta de la categoría de conjuntos (a saber, la categoría de preesquemas o la de preesquemas sobre un preesquema dado) nociones familiares para los conjuntos: productos cartesianos, leyes de grupo, de anillo, de módulo, fibrados, fibrados principales homogéneos, etc.
En el futuro, seguramente será difícil para el matemático eludir este nuevo esfuerzo de abstracción, quizá bastante pequeño, después de todo, en comparación con el que hicieron nuestros predecesores al familiarizarse con la teoría de conjuntos.

\bigskip
\begin{center}*\quad*\quad*\end{center}
\bigskip

Las referencias se darán según el sistema decimal; por ejemplo, en III, 4.9.3, III indica el capítulo, 4 el párrafo y 9 la sección del párrafo.
Dentro de un mismo capítulo, se omitirá la indicación del capítulo.
